\section{Research Objectives and Questions}

% The first objective (RQ1) investigates facilitating
% strategies and challenges in remote meetings to establish a foundation based on real-world practices and pain
% points. The second objective (RQ2) focuses on designing an AI-based moderation system that translates these
% insights into technical requirements, ensuring the system is context-aware and goal-oriented. The third
% objective (RQ3) examines how knowledge workers perceive a convergent remote meeting when a
% proactive AI moderator is present.

This thesis investigates support during a goal-oriented remote meeting, and it narrows that investigation
to the convergent phase: when a group must evaluate existing options, remain aligned with a
stated objective, and move toward a joint decision.
Within that scope, the work asks how a proactive, context-sensitive AI moderator can intervene
while the meeting is still running, especially where participants may not necessarily be aware that
something is not going ideally or is even counterproductive.

This thesis addresses the problem with a structured approach that integrates three complementary research
objectives. This bridges the gap between theoretical facilitation techniques and a practically applicable
AI-supported prototype. This prototype integrates a system proposal that builds upon existing HCI research
and concretizes it through features derived from domain-specific knowledge in the field of facilitation. The
result is a functional system that can be further evaluated and is intended to serve as a basis for further
iterations or as a source of inspiration for similar systems.

% Diese ARbeit macht dem nach eine großen abriss über die aktuelle forschung zu KI-gestützen Meeting Support
% sowie bestehende problme in remotmeetings. Um den fokus der arbeit einzuschreken leigt der fokus dabei auf
% der divergenten Phase von entscheidungsprozessen in MEetings, weil tools zur vor und nachbereitung berits
% teilweise in kommerziellen Meetinglösugnen zufinden sind und auch zu convergentn methoden ansätze
% besteheden,w ei KI unterstüzt (z.B. bei Brainstorming session etc.)

% ---------------------------------------------------------------------------------------------------------------------
\subsection{RQ1: Facilitating strategies and challenges in goal-oriented remote meetings}
\begin{quote}
  \textit{Which facilitating strategies are suitable for the goal-oriented management of remote meetings with
    convergent objectives (such as decision-making, prioritization, problem-solving)? What specific challenges
  typically arise in this context?}
\end{quote}
The first research question aims to identify effective strategies for managing convergent remote meetings,
such as decision-making, prioritization, and problem-solving, while also analyzing the specific challenges
that arise. To achieve this, the thesis conducts expert interviews with certified facilitators, agile
coaches, and team leads to gather empirical insights.
It also analyzes literature on facilitation techniques,
including timeboxing and role assignment, as well as the dynamics of remote collaboration. The findings are
compared with existing research to derive best practices and identify pain points, such as off-topic
discussions and unequal participation, which AI could potentially address.

% ---------------------------------------------------------------------------------------------------------------------
\subsection{RQ2: Design requirements for an AI-based moderation system}
\begin{quote}
  \textit{How should an AI-based moderation system be designed that autonomously detects thematic deviations
    based on a predefined meeting goal and real-time transcript analysis and promotes goal orientation through
  constructive, context-sensitive interventions?}
\end{quote}
The second research question focuses on developing a proactive AI system that detects thematic deviations in
real time and intervenes contextually to maintain goal orientation. This objective is addressed by defining
functional requirements, such as real-time transcript analysis and adaptive interventions, and non-functional
requirements, including privacy, latency, and scalability. A system architecture is then designed with
components for transcription, NLP-based analysis, and intervention logic, incorporating mechanisms for
private nudges or group reminders. A prototype is implemented using NLP libraries like Hugging Face
Transformers, and its technical feasibility is evaluated.

% ---------------------------------------------------------------------------------------------------------------------
\subsection{RQ3: Perception of the meeting with a proactive AI moderator}
\begin{quote}
  \textit{How do knowledge workers perceive a goal-oriented remote meeting in which a proactive AI
    moderator intervenes, with regard to communication, quality of discussion, status effects, teamwork,
  and efficiency, and how do they experience the agent's contributions?}
\end{quote}
The third research question investigates how knowledge workers experience a convergent remote
meeting that is supported by the prototype. The evaluation uses a hidden-profile hiring simulation
in which relevant information is distributed asymmetrically, so that a well-founded group decision
depends on sharing unique information. Meeting outcomes are measured with the questionnaire
instrument revalidated by Davison, covering communication, quality of discussion, status effects,
teamwork, and efficiency. Because that instrument was developed for earlier group support systems,
additional items address how participants experience the agent's interventions, including timing,
context-appropriateness, and whether feedback from the AI is more comfortable than equivalent
feedback from a group member.
